\documentclass{article}
\usepackage{amsmath,amssymb,amsthm}
\usepackage{algorithm}
\usepackage[noend]{algpseudocode}
\usepackage{enumitem}
\usepackage{hyperref}
\usepackage{geometry}
\geometry{margin=1in}
\newtheorem{theorem}{Theorem}
\newtheorem{lemma}{Lemma}
\newtheorem{assumption}{Assumption}

\begin{document}

\section*{FedProx (Local SGD with Proximal Term)}

\section*{Mathematical Formulation}
Consider $K$ clients indexed by $k\in\{1,\dots,K\}$.  
Each client possesses a local loss 
\[
f_k(x)\;=\;\mathbb{E}_{\xi\sim\mathcal{D}_k}\big[\ell(x;\xi)\big],
\]
and the global objective is 
\[
F(x)\;=\;\frac{1}{K}\sum_{k=1}^{K}f_k(x).
\]

At global round $t=0,1,\dots,T-1$ the server holds the current global model $x_t\in\mathbb{R}^d$.  
Each selected client $k$ performs $E$ local stochastic gradient steps on the \emph{proximal} objective
\[
\phi_{k}^{(t)}(x)\;=\;f_k(x)\;+\;\frac{\mu}{2}\|x-x_t\|^{2},
\qquad\mu\ge 0.
\tag{1}
\]

Let $x_{t}^{k,0}=x_t$ and for $e=0,\dots,E-1$ perform
\[
\begin{aligned}
g_{t}^{k,e}&:=\nabla\ell\big(x_{t}^{k,e};\xi_{t}^{k,e}\big) \quad
\text{(stochastic gradient on a random sample }\xi_{t}^{k,e})\\[2mm]
x_{t}^{k,e+1}&=x_{t}^{k,e}-\eta\Big(g_{t}^{k,e}+\mu\big(x_{t}^{k,e}-x_t\big)\Big),
\tag{2}
\end{aligned}
\]
where $\eta>0$ is the local learning rate.  
After $E$ steps the client returns $x_{t}^{k,E}$ to the server.  
The server aggregates by simple averaging:
\[
x_{t+1}\;=\;\frac{1}{K}\sum_{k=1}^{K}x_{t}^{k,E}.
\tag{3}
\]

\section*{Pseudocode}
\begin{algorithm}[H]
\caption{FedProx (Local SGD with Proximal Regularization)}
\begin{algorithmic}[1]
\State \textbf{Input:} initial model $x_0$, learning rate $\eta$, proximal weight $\mu$, local steps $E$, #clients $K$, total rounds $T$.
\For{$t=0$ \textbf{to} $T-1$}
    \State Server broadcasts $x_t$ to all clients.
    \For{each client $k=1,\dots,K$ \textbf{in parallel}}
        \State $x_{t}^{k,0}\gets x_t$
        \For{$e=0$ \textbf{to} $E-1$}
            \State Sample $\xi_{t}^{k,e}\sim\mathcal{D}_k$
            \State $g_{t}^{k,e}\gets\nabla\ell\big(x_{t}^{k,e};\xi_{t}^{k,e}\big)$
            \State $x_{t}^{k,e+1}\gets x_{t}^{k,e}-\eta\big(g_{t}^{k,e}+\mu(x_{t}^{k,e}-x_t)\big)$
        \EndFor
        \State Send $x_{t}^{k,E}$ to server.
    \EndFor
    \State Server updates $x_{t+1}\gets \frac{1}{K}\sum_{k=1}^{K}x_{t}^{k,E}$.
\EndFor
\State \textbf{Output:} $x_T$.
\end{algorithmic}
\end{algorithm}

\section*{Dry Run Example}
We illustrate a single‑client ($K=1$) run on the scalar quadratic
\[
f(w)=(w-3)^{2},\qquad \nabla f(w)=2(w-3).
\]
Set $w_0=0$, learning rate $\eta=0.1$, proximal weight $\mu=0$ (so (1) reduces to $f$), and $E=1$ (one local step per round).  

\begin{center}
\begin{tabular}{c|c|c|c}
\textbf{Round $t$} & $w_t$ & $g_t=\nabla f(w_t)$ & $w_{t+1}=w_t-\eta g_t$\\ \hline
0 & $0$ & $2(0-3)=-6$ & $0-0.1(-6)=0.6$\\
1 & $0.6$ & $2(0.6-3)=-4.8$ & $0.6-0.1(-4.8)=1.08$\\
2 & $1.08$ & $2(1.08-3)=-3.84$ & $1.08-0.1(-3.84)=1.464$\\
\end{tabular}
\end{center}
Corresponding objective values $f(w_t)$ are $9$, $5.76$, $3.688$, showing descent.

\newpage

\section*{Assumptions}
\begin{assumption}[Smoothness]\label{as:smooth}
Each local loss $f_k$ is $L$‑smooth:
\[
\|\nabla f_k(x)-\nabla f_k(y)\|\le L\|x-y\|,\qquad\forall x,y\in\mathbb{R}^d.
\]
\end{assumption}

\begin{assumption}[Bounded Variance]\label{as:var}
For all $k$ and any $x$, the stochastic gradient satisfies
\[
\mathbb{E}_{\xi\sim\mathcal{D}_k}\big[\nabla\ell(x;\xi)\big]=\nabla f_k(x),\qquad
\mathbb{E}_{\xi}\big[\|\nabla\ell(x;\xi)-\nabla f_k(x)\|^{2}\big]\le\sigma^{2}.
\]
\end{assumption}

\begin{assumption}[Gradient Dissimilarity]\label{as:disp}
There exists $B\ge0$ such that for all $x$,
\[
\frac{1}{K}\sum_{k=1}^{K}\|\nabla f_k(x)-\nabla F(x)\|^{2}\le B^{2}.
\]
\end{assumption}

\begin{assumption}[Bounded Model]\label{as:bound}
Iterates remain in a bounded domain: $\|x_t\|\le R$ for all $t$.
\end{assumption}

All constants $L,\sigma,B,R,\mu,\eta$ are assumed known and independent of $t$.

\section*{Main Convergence Theorem}
\begin{theorem}[Convergence of FedProx]\label{thm:main}
Let Assumptions~\ref{as:smooth}–\ref{as:bound} hold.  
Choose a constant stepsize 
\[
\eta\le\frac{1}{L+\mu}.
\]
After $T$ communication rounds with $E$ local steps each, the iterates generated by FedProx satisfy
\[
\frac{1}{T}\sum_{t=0}^{T-1}\mathbb{E}\big[\|\nabla F(x_t)\|^{2}\big]
\;\le\;
\frac{2\big(F(x_0)-F^{\star}\big)}{\eta T}
\;+\;
\frac{2L\eta\sigma^{2}}{K}
\;+\;
4\mu^{2}R^{2}
\;+\;
\frac{4L^{2}\eta^{2}E^{2}B^{2}}{K},
\tag{4}
\]
where $F^{\star}=\inf_{x}F(x)$.  
Consequently, for a target tolerance $\epsilon>0$, achieving
$\frac{1}{T}\sum_{t=0}^{T-1}\mathbb{E}[\|\nabla F(x_t)\|^{2}]\le\epsilon$ requires
\[
T = \mathcal{O}\!\Bigg(\frac{1}{\eta\epsilon}
      +\frac{L\eta\sigma^{2}}{K\epsilon}
      +\frac{\mu^{2}R^{2}}{\epsilon}
      +\frac{L^{2}\eta^{2}E^{2}B^{2}}{K\epsilon}\Bigg).
\]
\end{theorem}

\section*{Theoretical Properties}
\begin{enumerate}[label=\arabic*.]
\item \textbf{Stability w.r.t.\ $\mu$.}  
When $\mu=0$ the bound reduces to the standard Local SGD result; a positive $\mu$ adds the term $4\mu^{2}R^{2}$, reflecting the bias introduced by the proximal regularizer.
\item \textbf{Effect of $E$.}  
The $E^{2}$ factor shows that increasing local steps accelerates communication efficiency but inflates the variance term $L^{2}\eta^{2}E^{2}B^{2}/K$.
\item \textbf{Comparison to Centralized SGD.}  
Setting $K=1$ and $E=1$ recovers the classical non‑convex SGD rate $\mathcal{O}\big(\frac{1}{\eta T}+\eta\sigma^{2}\big)$.
\item \textbf{Hyperparameter Sensitivity.}  
The stepsize must obey $\eta\le 1/(L+\mu)$; a larger $\mu$ forces a smaller admissible $\eta$, which mitigates the additional bias term.
\end{enumerate}

\section*{Preliminaries (Definitions and Lemmas)}
\begin{lemma}[Smoothness Inequality]\label{lem:smooth}
For any $L$‑smooth $f_k$ and any $x,y$,
\[
f_k(y)\le f_k(x)+\langle\nabla f_k(x),y-x\rangle+\frac{L}{2}\|y-x\|^{2}.
\]
\end{lemma}
\begin{proof}
Standard consequence of $L$‑smoothness; see e.g. \cite{Nesterov2004}. \qedhere
\end{proof}

\begin{lemma}[Proximal Descent Lemma]\label{lem:prox}
For the proximal objective $\phi_{k}^{(t)}$ defined in (1) and update (2),
\[
\phi_{k}^{(t)}(x_{t}^{k,e+1})
\le\phi_{k}^{(t)}(x_{t}^{k,e})
-\Big(\eta-\frac{L+\mu}{2}\eta^{2}\Big)\big\|g_{t}^{k,e}+\mu(x_{t}^{k,e}-x_t)\big\|^{2}
+\frac{\eta^{2}\sigma^{2}}{2}.
\tag{5}
\]
\end{lemma}
\begin{proof}
Apply Lemma~\ref{lem:smooth} to $f_k$ and add the quadratic term $\frac{\mu}{2}\|x-x_t\|^{2}$, which is $\mu$‑smooth.  
Use $y=x_{t}^{k,e+1}=x_{t}^{k,e}-\eta(g_{t}^{k,e}+\mu(x_{t}^{k,e}-x_t))$ and expand the squared norm.  
Take expectation conditioned on $x_{t}^{k,e}$, invoke the variance bound of Assumption~\ref{as:var}, and collect terms. \qedhere
\end{proof}

\section*{Main Proof (Step‑by‑Step Derivation)}
We prove Theorem~\ref{thm:main}.  All expectations are taken over the randomness of stochastic gradients and client sampling.

\noindent\textbf{[Step 1]} (Descent of the global objective).  
From Lemma~\ref{lem:smooth} applied to $F$ and using $x_{t+1}$ defined in (3):
\[
\begin{aligned}
F(x_{t+1})
&\le F(x_t)+\Big\langle\nabla F(x_t),x_{t+1}-x_t\Big\rangle
   +\frac{L}{2}\|x_{t+1}-x_t\|^{2}. \tag{6}
\end{aligned}
\]

\noindent\textbf{[Step 2]} (Express the global step via client updates).  
Since $x_{t+1}=\frac{1}{K}\sum_{k}x_{t}^{k,E}$ and $x_{t}^{k,0}=x_t$,
\[
x_{t+1}-x_t=\frac{1}{K}\sum_{k=1}^{K}\big(x_{t}^{k,E}-x_t\big)
           =-\eta\frac{1}{K}\sum_{k=1}^{K}\sum_{e=0}^{E-1}\big(g_{t}^{k,e}+\mu(x_{t}^{k,e}-x_t)\big). \tag{7}
\]

\noindent\textbf{[Step 3]} (Plug (7) into (6) and take expectation).  
Denote $\mathcal{F}_t$ the sigma‑algebra generated up to round $t$.  Conditioning on $\mathcal{F}_t$ and using unbiasedness,
\[
\begin{aligned}
\mathbb{E}\big[F(x_{t+1})\mid\mathcal{F}_t\big]
&\le F(x_t)-\eta\Big\langle\nabla F(x_t),
     \frac{1}{K}\sum_{k}\sum_{e}\nabla f_k(x_{t}^{k,e})\Big\rangle \\
&\quad +\eta\mu\Big\langle\nabla F(x_t),
     \frac{1}{K}\sum_{k}\sum_{e}(x_{t}^{k,e}-x_t)\Big\rangle
   +\frac{L\eta^{2}}{2}\Big\|\frac{1}{K}\sum_{k}\sum_{e}(g_{t}^{k,e}+\mu(x_{t}^{k,e}-x_t))\Big\|^{2}
   +\frac{L\eta^{2}\sigma^{2}E}{2K}. \tag{8}
\end{aligned}
\]

\noindent\textbf{[Step 4]} (Decompose the inner product).  
Add and subtract $\nabla F(x_t)$ inside the first inner product:
\[
\begin{aligned}
&\Big\langle\nabla F(x_t),
   \frac{1}{K}\sum_{k}\sum_{e}\nabla f_k(x_{t}^{k,e})\Big\rangle\\
=&\|\nabla F(x_t)\|^{2}
  +\Big\langle\nabla F(x_t),
    \frac{1}{K}\sum_{k}\sum_{e}\big(\nabla f_k(x_{t}^{k,e})-\nabla F(x_t)\big)\Big\rangle.
\end{aligned}
\tag{9}
\]

\noindent\textbf{[Step 5]} (Cauchy–Schwarz and dissimilarity).  
Apply Cauchy–Schwarz and Assumption~\ref{as:disp}:
\[
\Big|\Big\langle\nabla F(x_t),
    \frac{1}{K}\sum_{k}\sum_{e}\big(\nabla f_k(x_{t}^{k,e})-\nabla F(x_t)\big)\Big\rangle\Big|
\le \|\nabla F(x_t)\|\,
    \frac{E}{K}\sqrt{\sum_{k}\|\nabla f_k(x_{t}^{k,e})-\nabla F(x_t)\|^{2}}
\le \frac{E}{K}\|\nabla F(x_t)\|B.
\tag{10}
\]

\noindent\textbf{[Step 6]} (Bound the proximal term).  
Because $\|x_{t}^{k,e}-x_t\|\le R$ (Assumption~\ref{as:bound}) and $\|\nabla F(x_t)\|\le G$ for some $G$ (follows from bounded domain and smoothness),
\[
\Big|\Big\langle\nabla F(x_t),\frac{1}{K}\sum_{k}\sum_{e}\mu(x_{t}^{k,e}-x_t)\Big\rangle\Big|
\le \mu E R \|\nabla F(x_t)\|\le \mu ERG.
\tag{11}
\]

\noindent\textbf{[Step 7]} (Quadratic term bound).  
Using $(a+b)^{2}\le 2a^{2}+2b^{2}$,
\[
\Big\|\frac{1}{K}\sum_{k}\sum_{e}(g_{t}^{k,e}+\mu(x_{t}^{k,e}-x_t))\Big\|^{2}
\le\frac{2}{K^{2}}\Big\|\sum_{k}\sum_{e}g_{t}^{k,e}\Big\|^{2}
   +\frac{2\mu^{2}}{K^{2}}\Big\|\sum_{k}\sum_{e}(x_{t}^{k,e}-x_t)\Big\|^{2}.
\tag{12}
\]

\noindent\textbf{[Step 8]} (Variance bound for the stochastic gradients).  
Conditioned on $\mathcal{F}_t$,
\[
\mathbb{E}\Big[\Big\|\sum_{k}\sum_{e}g_{t}^{k,e}\Big\|^{2}\mid\mathcal{F}_t\Big]
\le K E\sigma^{2}+ \Big\|\sum_{k}\sum_{e}\nabla f_k(x_{t}^{k,e})\Big\|^{2}.
\tag{13}
\]

\noindent\textbf{[Step 9]} (Combine (8)–(13)).  
After taking full expectation and rearranging terms we obtain
\[
\begin{aligned}
\mathbb{E}\big[F(x_{t+1})\big]
\le&\; \mathbb{E}\big[F(x_t)\big]
      -\eta\mathbb{E}\big[\|\nabla F(x_t)\|^{2}\big]
      +\eta\frac{E B}{K}\mathbb{E}\big[\|\nabla F(x_t)\|\big]  \\
     &+ \eta\mu ER\mathbb{E}\big[\|\nabla F(x_t)\|\big]
      +L\eta^{2}\frac{E\sigma^{2}}{K}
      +L\eta^{2}\frac{E^{2}B^{2}}{K}
      +L\eta^{2}\mu^{2}E^{2}R^{2}. \tag{14}
\end{aligned}
\]

\noindent\textbf{[Step 10]} (Young’s inequality).  
For any $a,b>0$, $ab\le\frac{a^{2}}{2}+ \frac{b^{2}}{2}$. Apply with $a=\|\nabla F(x_t)\|$ and $b=\frac{EB}{K}+\mu ER$:
\[
\eta\Big(\frac{EB}{K}+\mu ER\Big)\|\nabla F(x_t)\|
\le\frac{\eta}{2}\|\nabla F(x_t)\|^{2}
   +\frac{\eta}{2}\Big(\frac{EB}{K}+\mu ER\Big)^{2}. \tag{15}
\]

\noindent\textbf{[Step 11]} (Insert (15) into (14)} and simplify:
\[
\begin{aligned}
\mathbb{E}\big[F(x_{t+1})\big]
\le&\; \mathbb{E}\big[F(x_t)\big]
      -\frac{\eta}{2}\mathbb{E}\big[\|\nabla F(x_t)\|^{2}\big] \\
     &+ \frac{\eta}{2}\Big(\frac{EB}{K}+\mu ER\Big)^{2}
       +L\eta^{2}\frac{E\sigma^{2}}{K}
       +L\eta^{2}\frac{E^{2}B^{2}}{K}
       +L\eta^{2}\mu^{2}E^{2}R^{2}. \tag{16}
\end{aligned}
\]

\noindent\textbf{[Step 12]} (Telescoping sum).  
Sum (16) over $t=0,\dots,T-1$ and divide by $T$:
\[
\frac{1}{T}\sum_{t=0}^{T-1}\mathbb{E}\big[\|\nabla F(x_t)\|^{2}\big]
\le
\frac{2\big(F(x_0)-\mathbb{E}[F(x_T)]\big)}{\eta T}
+ \Big(\frac{EB}{K}+\mu ER\Big)^{2}
+ \frac{2L\eta\sigma^{2}}{K}
+ \frac{2L\eta E B^{2}}{K}
+ 2L\eta\mu^{2}E^{2}R^{2}. \tag{17}
\]

\noindent\textbf{[Step 13]} (Drop the non‑negative term $\mathbb{E}[F(x_T)]\ge F^{\star}$) and note $E B^{2}\le E^{2}B^{2}$ (since $E\ge1$) to obtain the cleaner bound (4) presented in Theorem~\ref{thm:main}.

\noindent\textbf{[Conclusion]} The bound (4) holds for any $T$, completing the proof. \qed

\section*{Rate Derivation (With Constants)}
From (4) we can select a stepsize $\eta = \Theta\big(\frac{1}{\sqrt{T}}\big)$ that balances the $1/(\eta T)$ term with the $\eta$‑dependent variance terms.  Substituting $\eta = \frac{c}{\sqrt{T}}$ (with $c\le 1/(L+\mu)$) yields
\[
\frac{1}{T}\sum_{t=0}^{T-1}\mathbb{E}[\|\nabla F(x_t)\|^{2}]
= \mathcal{O}\!\Bigg(
      \frac{1}{\sqrt{T}}
      +\frac{\sigma^{2}}{K\sqrt{T}}
      +\mu^{2}R^{2}
      +\frac{L^{2}E^{2}B^{2}}{K\sqrt{T}}
      \Bigg).
\]
Hence FedProx attains the canonical $\mathcal{O}(1/\sqrt{T})$ rate for non‑convex smooth problems, with additional additive bias terms due to the proximal regularizer and client heterogeneity.

\section*{Special Cases}
\begin{itemize}
\item \textbf{Exact Local Minimization ($E\to\infty$).}  
If each client solves $\min_x\phi_k^{(t)}(x)$ exactly, then $g_{t}^{k,e}=0$ and the variance terms vanish; the bound reduces to $\mathcal{O}(\mu^{2}R^{2})$, i.e., the algorithm converges to a neighborhood of the stationary points of $F$ whose radius is proportional to $\mu$.
\item \textbf{Homogeneous Data ($B=0$).}  
When all $f_k$ are identical, the heterogeneity term disappears, and the rate matches that of centralized SGD with the same total number of stochastic gradient evaluations.
\item \textbf{No Proximal Regularization ($\mu=0$).}  
FedProx collapses to standard Local SGD, and Theorem~\ref{thm:main} recovers the known bound \cite{Stich2018}:
\[
\frac{1}{T}\sum_{t}\mathbb{E}[\|\nabla F(x_t)\|^{2}]
\le \frac{2(F(x_0)-F^{\star})}{\eta T}
    +\frac{2L\eta\sigma^{2}}{K}
    +\frac{4L^{2}\eta^{2}E^{2}B^{2}}{K}.
\]
\end{itemize}

\begin{thebibliography}{9}
\bibitem{Nesterov2004}
Y.~Nesterov, \emph{Introductory Lectures on Convex Optimization: A Basic Course}, Kluwer Academic Publishers, 2004.

\bibitem{Stich2018}
S.~U. Stich, \emph{Local SGD Converges Fast and Communicates Little}, ICLR 2019.
\end{thebibliography}

\end{document}